\documentclass[10pt,a4paper]{amsart}
\usepackage[margin=19mm]{geometry}
\usepackage{amsmath,amssymb,mathtools}
\usepackage[scr=rsfs,cal=euler]{mathalfa}
\usepackage{xfrac}
\usepackage[unicode,hidelinks,bookmarks=false]{hyperref}
\newcommand{\ie}{i.e.\ }
\newcommand{\cf}{cf.\ }
\newcommand{\resp}{respectively\ }
\newcommand{\Subsection}{\subsection}
\providecommand{\nobreakdash}{\nobreak\hbox{-}}
\setlength{\emergencystretch}{2em}
\multlinegap0pt
\usepackage{bm}
\usepackage{bbm}
\usepackage{dsfont}
\usepackage{xspace}
\usepackage{cancel}
\usepackage{mathtools}
\usepackage{amsthm}
\makeatletter
\@ifundefined{theorem}{\newtheorem{theorem}{Theorem}}{}
\@ifundefined{lemma}{\newtheorem{lemma}[theorem]{Lemma}}{}
\makeatother
\title[On the Independence Numbers of the Cyclic Van der Waerden Hy]{On the Independence Numbers of the Cyclic Van der Waerden Hypergraphs}
\author{Benjamin Liber}
\date{Source: arXiv:2509.07926v2 (CC BY 4.0)}
\begin{document}
\maketitle

\noindent\emph{Source and licence.} On the Independence Numbers of the Cyclic Van der Waerden Hypergraphs. Version arXiv:2509.07926v2 (CC BY 4.0), distributed under the Creative Commons Attribution 4.0 International licence, as recorded in the arXiv licence metadata for this exact version and shown on the abstract page https://arxiv.org/abs/2509.07926v2. Full rights notice: licenses/arxiv-2509.07926-CC-BY-4.0.txt.\par\smallskip

\noindent\textit{Editorial scope.} Counted unit: the Lemma whose source label is \texttt{lem:cons} in the article (source0.tex lines 258--260, source label \texttt{lem:cons}), delivered together with the article's own complete original proof of it (source0.tex lines 262--294, one \texttt{proof} environment of 33 lines). No mathematics is written, repaired or re-derived: statement and proof are the article's own text line for line. Required context retained uncounted: lem:3.1real (lines 249--252). Every definition, display and label of those lines is retained unchanged. Omitted from this package and not redistributed: 1--248, 253--257, 261--261, 295--613. Nothing inside a retained excerpt is omitted or altered: every retained body is delivered exactly as the article's lines are written, so no omission receipt of any kind accompanies this package. Citations: the retained excerpts carry no citation command at all, so no bibliography entry of the article is transcribed and no reference is redistributed. Reference adaptation: none was needed and none was made; every reference inside the delivered unit resolves inside it, so no unresolved reference or citation is produced. Printed slips kept literal: no printed word, symbol, hypothesis or number of the counted statement or of its proof was altered. Layout and class: the article's own class is \texttt{article} with packages including graphicx, babel, geometry, amsmath, hyperref, amsthm, inputenc, amsfonts, enumerate, array, ulem, csquotes, xcolor, comment; the wrapper is the lane's pinned \texttt{amsart} editorial wrapper at 10pt on A4, which restates the theorem-like environments the retained text needs and adds the article's own macro definitions that the retained text needs (none), with extra wrapper packages (bm, bbm, dsfont, xspace, cancel, mathtools). The delivered numbering is the unit's own. Assets: no retained span contains a figure environment, a graphics-inclusion command, a PDF-page inclusion, a table or a TikZ picture; no image file is copied into this package, the package declares no asset dependency and no image file is redistributed with it. Licence: version arXiv:2509.07926v2 of this article is distributed under the Creative Commons Attribution 4.0 International licence, as recorded in the arXiv licence metadata for this exact version and shown on the abstract page https://arxiv.org/abs/2509.07926v2. A full rights notice is delivered at licenses/arxiv-2509.07926-CC-BY-4.0.txt. Characters: every retained excerpt is pure ASCII, so the delivered engine, XeLaTeX with its default Latin Modern text font, reports no missing character.
\begin{lemma}
\label{lem:3.1real}
    Let $d,N \in \mathbb{Z}^+$ such that $d \mid N$, and suppose $A$ is a cyclic arithmetic progression mod $N$ with common difference $d$. Then all residue classes of the elements of $A$ lie in a single congruence class of $\mathbb{Z}_{d}$.
\end{lemma}

\begin{lemma}\label{lem:cons}
Fix $d \in D(mk,k)$, and suppose $A$ is a $k$-term cyclic arithmetic progression mod $mk$ with common difference $d$. Let $\beta \in \{0,1,\ldots,d-1\}$ such that $A \subseteq [\beta]_d$\footnote{The existence of such a $\beta$ is guaranteed by Lemma \ref{lem:3.1real}.}. Set $\alpha := d-\beta$, and define \[S := \{k-\alpha,2k-\alpha,\ldots,mk-\alpha\} \subseteq \mathbb{Z}_{mk}.\] Then $A$ contains $d$ consecutive elements of $S$; that is, there exists an index $j \in \{1,\ldots,m\}$ such that \[\{jk-\alpha,(j+1)k-\alpha,\ldots,(j+d-1)k-\alpha\} \pmod {mk} \subseteq A,\] where the indices $j,j+1,\ldots,j+d-1$ are taken cyclically modulo $m$.
\end{lemma}

\begin{proof}
    First, since $d \in D(mk,k) = \{g \in \mathbb{Z}^+ : g \leq m, g \mid k\}$, we have that $d$ divides $k$. So, $k \equiv 0 \pmod d$. Hence, for every $j \in \{1,\ldots,m\}$, we have \begin{equation*}
        jk -\alpha \equiv jk-(d-\beta) \equiv -d+\beta \equiv \beta \pmod d.
    \end{equation*} So, $S \subseteq [\beta]_d$. On a similar note, since $A \subseteq [\beta]_d$, we can write $A = (a,a+d,a+2d,\ldots,a+(k-1)d) \pmod {mk}$ for some initial term $a \equiv \beta \pmod d$.

    Now, let us show that an arbitrary $x \in \mathbb{Z}_{mk}$ lies in $S$ if and only if $x \equiv -\alpha \pmod k$. First, suppose $x \in S$. Then $x \equiv jk-\alpha \pmod {mk}$ for some $j \in \{1,\ldots,m\}$. Reducing mod $k$ yields $x \equiv -\alpha \pmod k$. Conversely, suppose $x \equiv -\alpha \pmod k$. Then there exists an integer $q$ such that $x+\alpha = qk$. Take $j \in \{1,\ldots,m\}$ such that $j \equiv q \pmod m$ (if $q \equiv 0 \pmod m$, then take $j=m$). Writing $q = j+lm$ for some integer $l$, we have \begin{equation*}
        x+\alpha = qk = (j+lm)k = jk+l(mk).
\end{equation*} So, $x \equiv jk-\alpha \pmod {mk}$. Hence, $x \in S$. So, we see that $x \in S$ if and only if $x \equiv -\alpha \pmod k$.

Next, consider the following equation:\begin{equation}\label{equation1}
    dr \equiv -(\alpha+a) \pmod k.
\end{equation} Since $d$ divides $k$, we can write $k=td$ for some $t \in \mathbb{Z}^+$. So, we can rewrite (\ref{equation1}) as \begin{equation}\label{equation2}
    dr \equiv -(\alpha + a) \pmod {td}.
\end{equation} Since $a \equiv \beta \pmod d$ and $\alpha = d-\beta$, we have that \begin{equation*}
    \alpha+a \equiv \beta +(d-\beta) \equiv d \equiv 0 \pmod d.
\end{equation*} Hence, (\ref{equation2}) is solvable in terms of $r$. 

Now, suppose $r_0$ is one such solution to (\ref{equation2}). Then all solutions modulo $td$ are $r_0,r_0+t,r_0+2t,\ldots,r_0+(d-1)t$. Indeed, for each $j \in \{0,1,\ldots,d-1\}$, \begin{equation*}
    d(r+jt) \equiv dr+j(dt) \equiv dr \equiv -(\alpha+a) \pmod {td}.
\end{equation*} Moreover, suppose $r_0+lt \equiv r_0+l't \pmod {td}$ for some integers $0 \leq l,l' \leq d-1$. Then \begin{equation*}
    (l-l')t \equiv 0 \pmod {td} \implies l-l' \equiv 0 \pmod d.
\end{equation*} Since $|l-l'| < d$, it follows that $l-l' = 0$, i.e. $l=l'$. Hence, the solutions $r_0,r_0+t,r_0+2t,\ldots,r_0+(d-1)t$ are distinct modulo $td$. That is, there are $d$ distinct solutions to (\ref{equation2}) modulo $td$.

Finally, for such a solution $r_0$, consider the corresponding elements of $A$: \begin{equation*}
    A' :=\{a+dr_0,a+d(r_0+t),a+d(r_0+2t),\ldots,a+d(r_0+(d-1)t)\} \pmod {mk}.
\end{equation*} By construction, the first element of $A'$ satisfies \begin{equation*}
    a+dr_0 \equiv a-(\alpha+a) \equiv -\alpha \pmod k.
\end{equation*} So, $a+dr_0 \in S$. Hence, we can write $a+dr_0 \equiv jk-\alpha \pmod {mk}$ for some $j \in \{1,\ldots,m\}$. Each subsequent element in $A'$ differs from the previous one by $dt = k$. So, reducing $A'$ modulo $mk$ yields \begin{equation*}
    jk-\alpha, (j+1)k-\alpha,\ldots,(j+d-1)k-\alpha \pmod {mk},
\end{equation*} with indices taken cyclically modulo $m$. Hence, there exists an index $j \in \{1,\ldots,m\}$ such that \begin{equation*}
    \{jk-\alpha,(j+1)k-\alpha,\ldots,(j+d-1)k-\alpha\} \pmod {mk} \subseteq A,
\end{equation*} where the indices $j,j+1,\ldots,j+d-1$ are taken cyclically modulo $m$.
\end{proof}

\end{document}
